\subsection{Future Work}
\label{sec:discussion-future}

Several concrete directions follow from the present work. First, the training protocol relies on paired data from Ti-6Al-4V specimens; extending PMDF-Net to carbon fiber reinforced polymers and composite laminates will require domain adaptation techniques such as adversarial fine-tuning or physics-consistent transfer learning to account for material-dependent Lamb wave dispersion characteristics. Second, deployment in safety-critical inspection workflows demands confidence estimates alongside point predictions; Bayesian neural network layers or Monte Carlo dropout at inference can yield variance estimates that flag out-of-distribution inputs, enabling the network to abstain or trigger human review when uncertainty exceeds operational thresholds.

Third, real-time field deployment on portable platforms such as the Jetson Orin requires quantized network weights and profiling-driven operator fusion to meet latency constraints; systematic benchmarking on embedded hardware will establish the inference throughput achievable under representative signal batching. Fourth, the current architecture embeds physics constraints only through the loss function; hard-wiring dispersive Lamb wave dispersion relations as structured operations in the network topology could improve generalization to novel geometries without requiring geometry-specific retraining. Fifth, the paired data requirement limits applicability to domains where high-average reference signals are available; CycleGAN-based or adversarial denoising approaches can enable training from unpaired noisy/clean corpora, substantially expanding the reachable application space. Finally, the defect dataset covers primarily notches and flat-bottom holes; expanding to include fatigue cracks, stress corrosion cracking, and delamination would establish PMDF-Net performance across the full spectrum of structurally relevant flaw types encountered in industrial inspections.
